\section{School, Sisters, Rectory, Campus}

The parish's most consequential construction may have occurred outside the church. A national parish without a school could preserve worship in Polish while leaving the daily formation of children elsewhere. Saint Stanislaus began instruction very early and associated it with women religious, eventually the School Sisters of Notre Dame. Sources compete over whether the school began in 1866, 1867, 1872, or 1873, and whether a school building followed in 1873 or 1878. Those dates probably describe different things: informal or parish instruction, organization under a teaching community, and the first dedicated schoolhouse. The surviving evidence does not permit them to be forced into one event.

The institutional purpose is clearer than the disputed first day. Children learned religion and ordinary subjects while moving among Polish, German, and English. The parish's institutional history remembers rapid enrollment growth after religious sisters took charge and repeated enlargement of their residence and classrooms. The school made language transmission routine rather than ceremonial, trained later generations for civic and commercial life, and supplied the parish with choir members, societies, vocations, and family continuity.

Buildings followed enrollment. The earliest frame school gave way to a brick structure in the late nineteenth century. The present four-story school at 1669 South Fifth Street was designed by Herbst and Kuenzli; state preservation records date its plans and construction to 1926. Its Mediterranean or Spanish Revival façade fronted an older rear portion, which explains why the campus can look like one building while preserving multiple campaigns. A parish history's reference to a new school in the early 1930s may concern a later stage or educational expansion, but it cannot displace the dated plans without further archival evidence.

Secondary education expanded the parish's reach beyond its original grade-school role. Saint Stanislaus High School opened in September 1932 with seventy-two boys and girls, all freshmen. A 1946 \emph{Milwaukee Sentinel} report preserved by Marquette University counted four hundred students from thirty-seven parishes, taught by fourteen School Sisters of Notre Dame. The state inventory records the later name Notre Dame High School and a 1947 name change. The campus thus served not only neighborhood children but a wider network of Catholic families. Its classrooms were one more way the mother parish reproduced institutions outside its immediate boundaries.

Rectory and convent projects completed the working complex. An 1882 image caption identifies convent and rectory buildings, while later records show replacement or heavy alteration. In 1963--64 architect Mark Pfaller gave the rectory a new limestone façade and modern interior arrangement, and a connecting ambulatory and chapel joined it to the church. A later oral account called the older rectory razed and replaced; preservation records instead describe extensive remodeling and attachment. Until permit and property files are consulted, ``remade'' is safer than either extreme.

The campus should be read as an operating system. Church, school, convent, rectory, halls, and service spaces allowed clergy, sisters, children, families, societies, and the poor to encounter one another repeatedly. Its scale also created future vulnerability: when enrollment and neighborhood membership fell, the parish inherited far more masonry, roofing, heating, and maintenance than Sunday worship alone could sustain.

\evidenceframe{West Mitchell Street National Register nomination; Wisconsin AHI school record HI118139; ICKSP parish history; Notre Dame parish-history catalog.}{Instruction began early, religious sisters made it durable, and successive school and residential projects turned the parish into an educational campus.}{The precise inception of teaching and the first schoolhouse remain disputed. The 1926 plans securely control the present school's principal front campaign; a complete campus property history awaits parish, SSND, city-permit, and archdiocesan files.}
